\chapter{Introduction}
% This thesis argues that practical privacy is achieved not by assuming away
% existing infrastructure, but by modelling it accurately and using cryptography
% to constrain the information available to each participant according to their
% operational role.to start making the point that, though researchers interested
% in privacy usually use decentralisation to a very ... admirable? pov, a lot of
% the time in practice, solutions have a variety of parties and are not
% explicitly p2p, for example. and a way to successfully get privacy in practice
% is to embrace this.  think of privacy pass, apple private relay, ... apple hide
% my email, (finance stuff?) (physical cash itself?!) ... then looks at design
% docs for self sovereign identity, for ... digital assets? ??? and we see that
% to empower users is often through systems which allow them to achieve what they
% want in a frictionless, fast, secure, way, and this means that more powerful
% parties can be taken advantage of ...?  Privacy-enhancing technologies are
% increasingly deployed in systems that serve millions of users. These systems
% rarely consist of mutually distrustful peers. Instead, they involve service
% providers, issuers, auditors, validators, or other accountable entities that
% already have operational responsibilities.  Then examples.  Privacy Pass.
% Apple Private Relay.  Hide My Email.  Anonymous credential deployments.  CBDC
% proposals.  Enterprise blockchains.
%
% we dont (usually) trust them with privacy. they're honest but curious. they're
% trusted for availability, and to follow the protocol. and we can use this to
% build more lightweight systems.
%
% Real financial systems have legitimate accountability requirements.
%
% Modern privacy-preserving systems are rarely deployed in environments
% consisting solely of mutually distrusting peers. Instead, they are deployed
% within infrastructures that already contain organisations responsible for
% service provision, policy enforcement, auditing, or transaction validation.
% These organisations are not entrusted with users' privacy. Rather, they are
% assumed to execute prescribed protocols while cryptography limits the
% information available to them. This thesis explores how acknowledging these
% existing operational roles enables practical and efficient privacy-preserving
% systems.
%
% then the thesis,
%
% the aggregate credentials paper says, the credential systems we have at the
% moment don't reflect reality, so we make a nice neat construction that better
% reflects what is desired. and it's fast.
%
% the tracing paper, again, if we want financial privacy with digital money, we
% may need ability to both people at a transaction, and identify the user, and to
% point at a user, and identify all transactions associated with that user. this
% is what this paper enables. it uses epochs so that transactions aren't strictly
% sequential. i know tracing is a bad word for a lot of people. but maybe adding
% these features will allow us to have more privacy in digital money. (cite some
% stuff idk)
%
% offline payments. accessibility. develop the scheme, prove it secure, again
% there is an auditor. again talk about the power given to this person, contrast
% with solving for double spenders only, the efficiency / impossibility of that
%
% tsps. this is for a specific problem in hyperledger fabric style systems. there
% are known endorsers rather than untrusted anonymous miners, and so we want them
% to sign each transaction to confirm that it is legitimate. here is how you do
% that.
%
% This thesis argues that in practice, privacy is achieved not by assuming away
% existing infrastructure, but by modelling it accurately and using cryptography
% to constrain the information available to each participant according to their
% operational role.
